\subsection{CM theta normalization and further unconditional deductions}

The holomorphic theta normalization in Lead 16(e) differs from the
norm occurring in the question. We make that distinction explicit,
including the finite-character ambiguity in the multiplication law.

\begin{proposition}[ARI C.7.6; Proved: the exact CM multiplication law]\label{ari:C.7.6}
Let $E/\overline{\Q}$ have complex multiplication by an order
$\mathcal O$, and use an algebraic invariant differential. Write
$A=\operatorname{area}(\Lambda)$ and
\[
 \eta(z)=s_2z+\frac{\pi}{A}\bar z,\qquad
 \theta(z)=\sigma(z)e^{-s_2z^2/2}.
\]
Then $s_2\in\overline{\Q}$. For every nonzero
$\alpha\in\mathcal O$, with $d=\alpha\bar\alpha$, there is
$F_\alpha\in\overline{\Q}(E)$ such that
\begin{equation}\label{ari:eq:ct-CMmult}
 \theta(\alpha z)^2=F_\alpha(\exp_E z)\theta(z)^{2d}.
\end{equation}
Its divisor and leading Laurent coefficient are
\[
 \operatorname{div}(F_\alpha)=2[\alpha]^*(O)-2d(O),\qquad
 F_\alpha(\exp_Ez)=\alpha^2z^{2-2d}(1+O(z^2)).
\]
Thus the unsquared quotient is an algebraic function on $E$, but
it need not be a rational function on $E$. For a nontorsion
$P=\exp_Eu\in E(\overline{\Q})$, algebraicity of $\theta(u)$
would imply algebraicity of $\theta(\alpha u)$ for all
$\alpha\in\mathcal O$. The additive group $\mathcal O u$ has
rank two, not more than four.
\end{proposition}
\begin{proof}
The Legendre relation gives the displayed real-linear decomposition
of $\eta$. The sigma automorphy formula gives
\[
 \theta(z+\omega)=\epsilon(\omega)
  \exp\left(\frac{\pi\bar\omega}{A}
                         (z+\omega/2)\right)\theta(z),
 \qquad \epsilon(\omega)\in\{1,-1\}.
\]
For $\alpha\Lambda\subset\Lambda$, the exponential multipliers
of $\theta(\alpha z)/\theta(z)^d$ cancel because
$\alpha\bar\alpha=d$. Its remaining multiplier is a sign.
Squaring therefore gives a meromorphic function on $E$, with the
stated divisor. This divisor is defined over $\overline{\Q}$.
The existence of a rational function with this divisor over $\C$
implies its existence over $\overline{\Q}$: the relevant
one-dimensional Riemann--Roch space is obtained by extension of
scalars. Normalize its leading Laurent coefficient, in the formal
parameter of the algebraic differential, to $\alpha^2$. This
identifies it with the squared analytic quotient. The action of
$\alpha$ on this differential is algebraic, so $\alpha^2$ is an
allowed normalization.

Choose a nonreal $\alpha\in\mathcal O$. Since the normalized
sigma series is $\sigma(z)=z+O(z^5)$ with algebraic Laurent
coefficients, the next coefficient of the squared quotient is
\[
 \alpha^2z^{2-2d}
       \{1-s_2(\alpha^2-d)z^2+O(z^4)\}.
\]
The coefficients of the rational function in this algebraic formal
parameter are algebraic. As $\alpha^2-d\ne0$, this proves
$s_2\in\overline{\Q}$ without an unnormalized appeal to a CM
weight-two form. It also proves that the functions in
\eqref{ari:eq:ct-CMmult} are algebraic over the chosen field.
At nontorsion $u$, none of the arguments $\alpha u$ is a lattice
point; evaluating the identity proves the algebraicity implication.
The injective map $\mathcal O\to\C$, $\alpha\mapsto\alpha u$,
proves the rank assertion.

The qualification about rationality is necessary. For a square
lattice, $s_2=0$ and $\alpha=1+i$ give
$F_{1+i}=2i\wp$. The unsquared quotient has a simple zero at
the nonzero point in $\ker(1+i)$ and a simple pole at $O$,
which is not a principal divisor on $E$; its square is the
rational function just displayed.
\end{proof}

\begin{proposition}[ARI C.7.7; Known with derivation: two unconditional normalized values]\label{ari:C.7.7}
Let $E/\overline{\Q}$ have an algebraic invariant differential,
and let $u$ be an elliptic logarithm of an algebraic nontorsion
point $P$. Then
\[
 \sigma(u)\exp\left(-\frac12u\zeta(u)\right)
\]
is transcendental. If $E$ has CM and $s_2$ is normalized as in
C.7.6, then $\zeta(u)-s_2u$ is also transcendental. The first
assertion concerns the displayed holomorphic normalization;
it does not assert transcendence of either $N(P)$ or
$\theta(u)=\sigma(u)e^{-s_2u^2/2}$.
\end{proposition}
\begin{proof}
We use \citet[Theorem 1 and Corollary 1]{waldschmidt1979sigma},
whose complete proof was opened in the translation hosted by its
author. Its Theorem 1 states that, for algebraic invariants
$g_2,g_3$, algebraic elliptic points $\exp_Eu$, $\exp_Ev$, and
$\beta\in\overline{\Q}$, the number
\[
 \frac{\sigma(u+v)}{\sigma(u)\sigma(v)}
                      \exp((\beta-\zeta(v))u)
\]
is transcendental if $\exp_Ev$ is nontorsion and $u,u+v$
are not lattice points. Here ``algebraic elliptic point''
means that the Weierstrass coordinates are algebraic; it does
not mean that its logarithm is algebraic. Take $v=u$ and
$\beta=0$. All hypotheses hold because $P$ is nontorsion.
The sigma duplication formula identifies the displayed number
with
\[
 -\wp'(u)\sigma(u)^2e^{-u\zeta(u)}.
\]
The factor $-\wp'(u)$ is nonzero and algebraic. Therefore
$\sigma(u)e^{-u\zeta(u)/2}$ is transcendental, as also stated
explicitly after Corollary 1 of the cited paper.

For the second assertion we apply the analytic subgroup theorem
in \citet[Theorem 6.1]{huber2022}. In the algebraic universal
vector extension $\widetilde E$ a logarithm of an algebraic lift
of $P$ has coordinates $(u,\zeta(u)-b)$, where $b$ is an algebraic
fibre coordinate, with a possible algebraic linear change of the
second coordinate. If $b=\zeta(u)-s_2u$ were algebraic, this
logarithm would belong to the algebraic line
$\{(z,s_2z):z\in\C\}$. The theorem applies to the connected
commutative algebraic group $\widetilde E$ over
$\overline{\Q}$ and this line, defined over $\overline{\Q}$ by
C.7.6; its exponential contains a nonzero algebraic lift.
It would provide a connected one-dimensional algebraic subgroup
whose Lie algebra is that line. Its projection to $E$ is
nontrivial. A one-dimensional affine algebraic group has no
nonconstant algebraic-group homomorphism to an elliptic curve,
so the subgroup is an elliptic curve isogenous to $E$.
Its inclusion would split the universal vector extension after
that isogeny. This is impossible: the nonzero extension class in
$\operatorname{Ext}^1(E,\mathbb G_a)=H^1(E,\mathcal O_E)$
remains nonzero under isogeny, since pullback is invertible in
characteristic zero.
\end{proof}

\paragraph{A theta criterion requiring its original normalization.}
An earlier draft applied the pointwise criterion printed in
\citet[final slide, comment (c)]{bertrand2017} to
$g(u,v)-s_2uv$, where
$g(u,v)=\log(\sigma(u+v)/(\sigma(u)\sigma(v)))$.
It inferred transcendence of $\theta(u)$ by taking $v=u$.
We withdraw that inference pending verification of the original
lift conventions in \citet[Corollary 3]{bertrand1984}.
Indeed replacing $v$ by $v+\omega$, $\omega\in\Lambda$,
changes $g(u,v)-s_2uv$, modulo $2\pi i\Z$, by
\[
 (\eta(\omega)-s_2\omega)u=\frac{\pi\bar\omega u}{A},
\]
whereas the condition $N\exp_Ev=f(\exp_Eu)$ with $f$
Rosati-antisymmetric depends only on the elliptic points.
Thus the two-sided statement cannot be used with unrestricted
lifts without checking its normalization. The author's slide
and the original paper's first two pages were accessible;
the full original Corollary 3 was not. We use neither direction
of that criterion here. The withdrawn text and the lift audit
are retained in the records.

A subsequent source check resolves the normalization of the
one-variable holomorphic assertion itself:
\citet[p.~15-03]{waldschmidt1988theta} explicitly states that
$\sigma(u)e^{-s_2u^2/2}$ is transcendental when the invariants
are algebraic, $E$ has CM, and $\exp_Eu$ is an algebraic
nontorsion point.  That source defines $s_2$ by the regularized
lattice sum and attributes the result to Bertrand's 1982
paper \emph{Valeurs de fonctions th\^eta et hauteurs p-adiques}.
We have checked this complete normalized statement in
Waldschmidt's published exposition; the full original 1982
proof has not been retrieved.  This source clarification
does not validate the unrestricted-lift inference withdrawn
above, and the theorem concerns the holomorphic value rather
than the metric norm.  C.7.7 retains the two assertions whose
proofs have been fully checked here.

\paragraph{The archimedean term still requiring control.}
The exact comparison is
\begin{equation}\label{ari:eq:ct-Gaussian}
 N(P)=|\Delta|^{1/12}|\theta(u)|
                         \exp\left(-\frac{\pi|u|^2}{2A}\right).
\end{equation}
It follows from C.7.1 and the real-linear decomposition of
$\eta$. The normalization proved transcendental in C.7.7 instead
gives the exact identity
\[
 N(P)=|\Delta|^{1/12}
 \left|\sigma(u)e^{-u\zeta(u)/2}\right|
 \exp\left(\frac12\Re\{u(\zeta(u)-\eta(u))\}\right).
\]
Neither formula excludes algebraicity of $N(P)$: transcendence
of one complex factor does not control its absolute value or
its product with the remaining exponential.
The generic finite-order interpolation approach in Lead 16(e)
supplies only the rank-two set $\mathcal O u$; it does not
meet the proposed rank-greater-than-four criterion. The
derivative approach would need the algebraicity of
$\theta'(u)/\theta(u)=\zeta(u)-s_2u$, which C.7.7 excludes.
Replacing the norm by the holomorphic theta value would
therefore change the problem precisely by the factor in
\eqref{ari:eq:ct-Gaussian}.

The related independence theorem of
\citet[Corollary~2]{reyssat1982weierstrass} says that two of
$e^{i\pi u/\omega_1}$, $s(u)$ and $s'(u)$ are algebraically
independent, where $s(z)=\sigma(z)e^{-\eta_1z^2/(2\omega_1)}$.
For real $u$, hypothetical algebraicity of $N(P)$ makes
$s(u)$ algebraic, so the conclusion instead forces independence
of $e^{i\pi u/\omega_1}$ and
$\zeta(u)-\eta_1u/\omega_1$; it yields no contradiction.
The periodic-functional-equation theorem reported in
\citet[p.~15-05]{waldschmidt1988theta} assumes algebraic
$q=e^{i\pi\tau}$, whereas this $q$ is transcendental for
imaginary quadratic $\tau$ by Gelfond--Schneider.  These
specific checked applications therefore do not decide the
metric norm; their source scopes and the larger CM
independence tuples are recorded separately.

We also tested an additive-extension route.
\citet[Theorem 3.1.2]{bost2009} proves that an algebraic point
of the universal vector extension of an abelian variety which
lies in its maximal compact subgroup at one complex embedding
is torsion. The theorem's hypotheses require an algebraic
point of the vector extension, equivalently an algebraic line
bundle with an algebraic connection having unitary monodromy.
C.7.2 supplies only an algebraic point in one fibre of a
multiplicative extension. It does not supply that algebraic
connection. Differentiating a metric value at one algebraic
point cannot establish algebraicity of the connection there
or globally. This is the exact missing implication in that
attempt.

Another route uses the Latt\`es map $f_E$ defined by
$x(2Q)=f_E(x(Q))$. Its canonical height is
$\widehat h_{f_E}(x(P))=2\widehat h_O(P)$. For rational
parameters C.6 therefore makes algebraicity of $N(P)$
equivalent to algebraicity of
$\exp(\widehat h_{f_E}(x(P)))$.
\citet[Theorem 1.4 and Corollary 1.6]{nguyen2023} settle the
corresponding multiplicative-height problem for nonintegrable
polynomials using B\"ottcher coordinates in the basin of
infinity. They do not apply to $f_E$. Indeed an elliptic
duplication Latt\`es map has no exceptional point: the inverse
images of any point contain the images of all successive
division points of an elliptic point, an infinite set. A
polynomial has the exceptional point $\infty$, so no conjugate
of $f_E$ is polynomial. The Julia set is the whole sphere
and the required basin of infinity is absent. This checks
the failed hypothesis, rather than treating a theorem for
polynomials as a theorem for all rational maps.

We also checked the arithmetic algebraicity criterion of
\citet[Theorems 1.3 and 2.3]{xie2024algebraicity}.
Besides sufficiently rapid height-normalized approximation by
points in a fixed number field, it requires the analytic germ
to be formally defined over that field. This hypothesis fails
for the compact leaf in C.7.9: its tangent has the proved
transcendental slope
$c=\zeta(u)-\eta_1u/\omega_1$.
A smooth formal germ defined over $\overline{\Q}$ has an
algebraic tangent; algebraic translations, coordinate changes
and isogenies preserve this obstruction. Assuming $N(P)$
algebraic supplies an algebraic point on the leaf and its
integer multiples, but not an algebraic formal germ. Attempting
to force convergence by successive division also loses the
fixed-number-field hypothesis: if $[2]^nQ_n=Q$ with all
$Q_n\in E(K)$ for one number field, then
$\widehat h_O(Q_n)=4^{-n}\widehat h_O(Q)>0$.
The $Q_n$ are distinct with bounded ample height, contrary to
Northcott finiteness. The full hypothesis audit and the opened
preprint are retained in the records.

\paragraph{A current source whose theorem text remains unchecked.}
The publisher's abstract for \citet{pham2025sigma} announces
complex and $p$-adic linear-independence results for semiabelian
varieties and applications to Weierstrass sigma functions.
We opened its metadata, abstract and reference list, but both
the full-HTML link and the publisher's citation-PDF link returned
the abstract/access page. No accessible author manuscript was
located. Consequently we could not check any theorem statement
or its hypotheses. This is a potentially relevant unverified
source: we make no claim either that it settles the Gaussian
norm question or that its results are inapplicable to it.

